\section{Contemporary Requirements for Magical Idealism}

\begin{quotex}
\emph{Evola invests the dark matter of esoterism with a dialectic and discursive reflexion}. 
\flright{From Franco Volpi's introduction to Magical Idealism}
\end{quotex}


\begin{quotationx}
TRANSLATION of Chapter 6 of \emph{Essays on Magical Idealism} (\emph{Saggi sull'idealismo magico}) by \textbf{Julius Evola}.

In this early book, the young Evola sketched out his philosophical position which he named ``Magical Idealism'' (presumably inspired by \textbf{Novalis}). Although this predates his interest in Tradition aroused in him by Rene Guenon, he never repudiated his philosophy, but rather tried to integrate it with his understanding of Tradition.

It is easy to forget in our time how influential philosophical idealism was in the 1920s. It was still dominant in Germany after more than a century of development. England had Green, Bradley, Bosanquet, and Collingwood. France, too, had its own stream which reached its peak in the spiritualism of Lavelle and Lesenne (both indirectly influenced by Octave Hamelin.) And of course, Italian philosophy was dominated by the idealists Benedetto Croce and Giovanni Gentile. Evola sought to insinuate himself in that current of thought. Actually, as this introduction shows, he regarded his own system as the culmination and fulfillment of what went before.

Evola includes brief summaries of five thinkers who most influenced magical idealism. These are \textbf{Carlo Michelstaedter}, \textbf{Otto Braun}, \textbf{Giovanni Gentile}, \textbf{Octave Hamelin}, and \textbf{Herman Keyserling}. Of these, I will be offering translations of the first, third, and the fifth, since I am familiar with them and have read their works. I know nothing of Braun and Hamelin, both very obscure today, so I doubt I will bother to translate those sections.
\end{quotationx}


A requirement raised many times in these pages is to show how the conception of magical idealism carries on by way of logical continuity and integrates the most advanced positions that modern Western speculative thinking has achieved. We now want to give a specific, direct satisfaction to one such requirement by considering a group of thinkers among the most significant in contemporary culture, emphasizing the deep theme that informs their conceptions and finally showing how, when such a theme is given free reign in the interior of their own systems --- and without violating any of their existing parts, but rather bringing them to a great organic perfection --- we reach those assertions that were sketched out in the preceding chapters.

However, it is first necessary to understand well the meaning of this \emph{historical conclusion} of magical idealism. It is that if we conceive history as existing in itself, as therefore imposing the bad misfortune of a group of \emph{given} elements from which, in one way or another, the current moment would come to be conditioned, a demonstration of the historical necessity of magical idealism in truth could have value only as a true rebuttal of magical idealism itself, since the fundamental principle of this doctrine is absolute, unconditioned self-determination. What cannot therefore happen if something stands against the I, which is simply given to it, something that is there without the participation of its will. Things proceed however quite differently when one holds firm to the principle of the ideality of time and, with it, of history. If time is not a thing in itself, but rather---as Kant taught---a category, if it is simply a mode by which the I orders the matter of representation which therefore, in itself, is neither temporal nor intemporal, it does not exist as a before nor as an after---then the spectre of an inevitable determination from part of the past vanishes into nothing. Since in such a case it instead remains true that, insofar as the past exists only inside the act---which in itself, in this regard, is to be understood as metatemporal---by which I make my various affirmations appear temporally, the past does not condition or determine the present, but the present conditions or determines the past.

The past remains simply a mark with which I identify a part of my current experience, since a past \emph{in itself}, i.e., a past that falls outside my real experience, which is not an object, is gnoseologically an absurdity and non-being. From that, it follows that history is nothing other than a mode according to which the I projects onto the canvas of time, I would say almost as in a mythical figuration, that it is to will interiorly and intemporally. As creator of history in the current historical moment, the individual experiences in this way only the limit-point of his own affirmation. The theory of the ideality of time therefore makes history a plastic faculty and in itself is indifferent to freedom---no longer a tyrannical fact that does violence to the individual, but rather a docile creation that the individual dominates, reflects and confirms unfailingly \emph{a posteriori} what it \emph{a priori} and metahistorically affirms. Thus, properly, it should be said that history is nothing other than the very power of freedom in reflecting and demonstrating \emph{a posteriori}, alongside the category of time, its determination having happened \emph{a priori} in an intemporal and metahistorical point. The ``historical conclusion'' is always something that comes after, an \emph{epiginomenon}, and its necessity is only the phenomenon of freedom that determines it unconditionally.

That granted, we can show the historical necessity of magical idealism without that assumption implying a contradiction. It rises like the synthesis of a dialecticism, in which the thesis is the rationalism of romantic philosophy which, exhausting itself in a conceptual world abstracted from reality and from individuality, generated the antithesis of materialism and positivism. 1 For the consummation of the thesis in the antithesis, the empty ideality had to be refilled with a concrete content whence, in the terms of the Hegelian Left (Stirner, Nietzsche), the mentioned emergence in the affirmation of the real individual in the value of the unconditioned. This principle of synthesis was developed, then it led to the concept of an individual affirmation such that in the same plane of the real world presented antithetically by positive science and established that sufficiency in that mediation, of which in the abstract world of the rational only the image without life was known. As the rationalistic thesis culminated in an idealization of the real, so a realization of the ideal was postulated of the synthesis of magical idealism (which is then the true derealization of the real), that is, a power of the individual so real, that he was the being and determination of nature studied from the antithetical moment of science. 1 The surpassing of the Hegelian Logos, so that this acquires its concreteness, is not that of nature, internal to the pure logical sphere of the Encyclopedia of Philosophic Sciences, but rather that of the entire Hegelianism, through the Hegelian left, in the sciences of nature, with the organs proper to which and falling outside the pure conceptual apriorism and took specific knowledge of concrete reality.


\flrightit{Posted on 2014-05-23 by Cologero}
